
\subsection{Consistency checks}\label{subsec: checks}

\paragraph{(a) No foreign rate.} If $\sigma^f_r=0$ then $\zeta^f=\zeta^d=\ell^f\equiv 0$ and the first line of (\ref{eq: sensitivities}) is (\ref{eq: oneProcess}): formula (\ref{eq: mainFormula}) is the one-factor Hull-White formula (\ref{eq: malliavinHullWhite}) of \cite{guyonLabordere}.

\paragraph{(b) Constant spot volatility.} If $\sigma(t,S)\equiv\sigma$ then $g=0$, the system (\ref{eq: sensitivities}) is deterministic and can be integrated:
$$ 2\big(\ell^d_t-\ell^f_t\big)=\partial_t\int_0^t\Big[
(\sigma^d_r A^d)^2+(\sigma^f_r A^f)^2+2\rho_{Sd}\sigma\sigma^d_r A^d-2\rho_{Sf}\sigma\sigma^f_r A^f
-2\rho_{df}\sigma^d_r\sigma^f_r A^dA^f\Big](s,t)\,ds . $$
The integral is the part of the total variance of the forward $f$ that is due to the rates in the Gaussian three-factor model, so the right hand side is exactly $\sigma_{Dup}(t)^2-\sigma^2$. With zero correlations this reads
$$ \sigma^2=\sigma_{Dup}(t)^2-\frac{(\sigma^d_r)^2}{\kappa_d^2}(1-e^{-\kappa_dt})^2-\frac{(\sigma^f_r)^2}{\kappa_f^2}(1-e^{-\kappa_ft})^2 . $$
This fixes all signs and the normalisation $2/K$ of both terms in (\ref{eq: threeFactorGeneralMalliavin}).

\paragraph{(c) Monte Carlo, with skew.} Integrating (\ref{eq: mainFormula}) twice in the strike gives two identities between unconditional expectations, which can be tested without any kernel estimate: with $c^i_T=r^i_T-r^0_i(T)$,
\begin{align*}
\mathrm{(I)}\quad&\mathbb{E}^d\big[D^d_{0T}\,c^d_T\,(S_T-K)^+\big]=\mathbb{E}^d\big[D^d_{0T}\,\mathbb{I}_{S_T>K}\,S_T\,\ell^d_T\big],\\
\mathrm{(II)}\quad&\mathbb{E}^d\big[D^d_{0T}\,S_T\,c^f_T\,(S_T-K)^+\big]=\mathbb{E}^d\big[D^d_{0T}\,\mathbb{I}_{S_T>K}\,S_T^2\,\ell^f_T\big].
\end{align*}
They were tested for a fixed local volatility with a pronounced skew,
$$\sigma(S)=0.12\,\big(1-0.6\tanh(2\ln S)\big),\qquad S_0=1,\qquad T=5,$$
with $\kappa_d=5\%$, $\sigma^d_r=1.2\%$, $\kappa_f=10\%$, $\sigma^f_r=1.5\%$, correlations
$\rho_{Sd}=0.35$, $\rho_{Sf}=-0.30$, $\rho_{df}=0.40$, flat curves at $3\%$ and
$1\%$, $2\cdot10^6$ antithetic paths and 200 Euler steps per year. In the table, $z$ is the difference between the two sides in standard errors. The column ``naive'' uses the system (\ref{eq: sensitivities}) without the cross terms and with $g(dW^S-\sigma dt)$ in $\ell^f$, i.e. without the corrections due to the measure change.

\begin{center}
\begin{tabular}{c|rrr|rrr|rr}
\toprule
 & \multicolumn{3}{c|}{(I) domestic, $\times10^{-3}$} & \multicolumn{5}{c}{(II) foreign, $\times10^{-3}$}\\
$K$ & direct & (\ref{eq: sensitivities}) & $z$ & direct & (\ref{eq: sensitivities}) & $z$ & naive & $z$\\
\midrule
0.70 & $2.425$ & $2.422$ & $+0.4$ & $-3.096$ & $-3.094$ & $-0.2$ & $-3.180$ & $+7.0$\\
0.85 & $2.127$ & $2.124$ & $+0.5$ & $-2.856$ & $-2.854$ & $-0.2$ & $-2.930$ & $+7.2$\\
1.00 & $1.712$ & $1.709$ & $+0.8$ & $-2.454$ & $-2.451$ & $-0.3$ & $-2.512$ & $+7.0$\\
1.15 & $1.236$ & $1.236$ & $0.0$ & $-1.914$ & $-1.914$ & $0.0$ & $-1.956$ & $+6.7$\\
1.30 & $0.787$ & $0.788$ & $-0.4$ & $-1.324$ & $-1.325$ & $+0.1$ & $-1.350$ & $+5.5$\\
1.50 & $0.349$ & $0.350$ & $-0.5$ & $-0.658$ & $-0.658$ & $+0.2$ & $-0.668$ & $+3.7$\\
\bottomrule
\end{tabular}
\end{center}

\noindent Both identities hold within one standard error, while the naive foreign term is off by about $2.7\%$. For these parameters the cross term in the domestic pair is small (about $0.2\%$, not resolved by the test). In a stressed configuration ($T=10$, both rate volatilities doubled, skew coefficient $0.8$, 100 steps per year) both identities again hold within one standard error, the naive domestic term is off by about $1\%$ ($z$ between $3$ and $5$) and the naive foreign term by about $3\%$ ($z$ between $4$ and $8$).

The simulation scheme of Section \ref{subsec: discretisation3F} was checked separately: the closed forms of all step coefficients, including $K^{df}_j$, agree with numerical quadrature to $10^{-10}$; under $\mathbb{Q}^{d^T}$ the simulated particles reproduce $\mathbb{E}[f_t]=f_0$, $P^d_{0T}\mathbb{E}[(P^d_{tT})^{-1}]=P^d_{0t}$ and $P^d_{0T}\mathbb{E}[(P^d_{tT})^{-1}S_t]=S_0P^f_{0t}$ within Monte Carlo error; and (I)--(II) hold at an intermediate date ($t=4$, $T=6$) with $\mathbb{Q}^{d^T}$ particles and weights $(P^d_{tT})^{-1}$. Pooling five independent runs of $10^6$ paths at that date (two under $\mathbb{Q}^{d^T}$, three under $\mathbb{Q}^d$), the two sides differ at $K=0.7$ by $0.13\%$ for (I) and $0.09\%$ for (II), that is $0.8$ and $0.4$ standard errors; single runs deviate by up to $2.9$ standard errors at one strike, and the right hand sides computed under the two measures agree to about $0.1\%$.
The matrix form (\ref{eq: matrixSystem}), simulated directly, agrees with (\ref{eq: sensitivities}) path by path up to the Euler error. Scripts: \texttt{check\_sensitivities.py}, \texttt{check\_forward\_measure.py}, \texttt{check\_matrix\_form.py}.

\subsection{Scope and caveats}\label{subsec: caveats}

\paragraph{Calibration identity versus existence of a calibrated model.} Propositions \ref{prop: calibration} and \ref{prop: main} are statements about a model that exists: \emph{if} (\ref{eq: threeFactorModel}) has a sufficiently regular solution, then it is calibrated if and only if its local volatility satisfies (\ref{eq: mainFormula}). Substituting (\ref{eq: mainFormula}) into (\ref{eq: threeFactorModel}) gives an SDE which is nonlinear in the sense of McKean, whose coefficient depends on the law of the solution through a conditional expectation given $S_t$. Neither existence nor uniqueness of a solution of this equation follows from the identity. Even for local-stochastic volatility models with deterministic rates, existence is known only in special cases: perturbative and local in time \cite{abergelTachet}, for regime-switching volatility under a condition on the range of the volatility factor \cite{jourdainZhou}, and for stationary solutions of a time-homogeneous version \cite{lackerShkolnikovZhang}. We are not aware of an existence result for the stochastic-rates equation considered here. The particle algorithm is therefore a numerical construction; its convergence as $N\to\infty$ is not proved in this paper.

\paragraph{Non-negativity of the local variance.} The right hand side of (\ref{eq: mainFormula}) need not be positive. Check (b) makes this explicit: with a flat smile, the calibrated variance is $\sigma_{Dup}(t)^2$ minus the forward variance generated by the rates, which is negative as soon as the rate volatilities and correlations imply more forward variance than the market quotes, typically at long maturities. In that case no local volatility model with the chosen rate parameters reproduces the smile; conditions under which a local volatility exists are discussed in \cite{ogetbilGanesanHientzsch}. In the algorithm, $\sigma_N(t,K)^2$ should be floored at a small positive value and the region of $(t,K)$ where the floor binds should be reported: there the model is not calibrated, and the remedy is in the inputs $(\kappa_i,\sigma^i_r,\rho)$, not in the numerics.

\paragraph{Regularity.} The derivation assumes that $\sigma$ is continuously differentiable in $S$ with $S\partial_S\sigma$ bounded (first variation process), that $S_t$ has a smooth positive density, that the stochastic exponentials defining the measure changes are true martingales, and enough integrability for the Malliavin integration by parts. In the algorithm, $g_j$ is the derivative of an estimated function: noise in $\partial_S\sigma_N$ enters $\ell^d,\ell^f$ multiplicatively. The same issue is present in the one-factor method through $V_t$.

\paragraph{What has and has not been tested.} The Monte Carlo evidence of Section \ref{subsec: checks} concerns the identities behind (\ref{eq: mainFormula}) for a \emph{given} smooth local volatility, and the simulation scheme. The full particle calibration of the three-factor model against a market smile, followed by the repricing of vanilla options, has not been run yet; nor has the bias due to the time step and to the kernel bandwidth been quantified.

\paragraph{Hull-White.} Proposition \ref{prop: main} uses two properties of the Hull-White model: the Clark-Ocone integrands (\ref{eq: kernels}) are deterministic and factorise, and the bond volatilities are deterministic, so that risk-neutral and forward measures differ by deterministic drifts.

\paragraph{Related literature.} We have not checked systematically whether the representation (\ref{eq: mainFormula}) appears in the literature cited at the beginning of this section; no claim of novelty is made at this stage.

\subsection{Outlook}
\paragraph{Gaussian multi-factor rates.} If a rate is a sum of Hull-White factors, the kernel (\ref{eq: kernels}) is a sum of exponentials and each factor contributes its own pair $(\ell,\zeta)$.

\paragraph{General one-factor short rates.} For rates $dr^i_t=\mu^i(t,r^i_t)dt+\sigma^i_r(t,r^i_t)dB^i_t$ the statement that replaces the factorisation $V_t/V_s$ is $\MD_sX_t=J_tJ_s^{-1}\mathcal{A}(s,X_s)\mathcal{C}^\top e$, with $J$ the $3\times3$ first variation matrix of $X=(S,r^d,r^f)$, which is not triangular because of the quanto drift. Two further difficulties appear: the Clark-Ocone integrand is stochastic and depends on $t$, and the forward measures differ from the risk-neutral ones by state-dependent drifts, which produce additional terms of the type of $H_t$.

\paragraph{Stochastic volatility.} With a stochastic volatility factor $a_t$ as in Section 1, the left hand side of (\ref{eq: mainFormula}) becomes $\sigma(t,K)^2\,\mathbb{E}^{d^t}[a_t^2|S_t=K]$ and the sensitivities acquire the tangent process of $a_t$ when $a_t$ is correlated with the rates. This is not worked out here.

\appendix
\section{Changes with respect to the 2016 draft}
\begin{itemize}
\item Section 1: $r^0_t$ is the forward rate $\mathbb{E}^{\mathbb{Q}^t}[r_t]$ (it was written $\mathbb{E}^{\mathbb{Q}}[r_t]$).
\item Section 2: exponent in $\theta(t)$; definition of $B^T_t$; $\Xi_t$ is defined with $\sigma_r$; missing $+\sigma_r$ in front of the stochastic integral in (\ref{eq: discretization_forwardPrice}); initialisation of the algorithm. New Section \ref{subsec: oneProcess}.
\item Spot SDE of the three-factor model: the diffusion term is $\sigma(t,S_t)S_tdW^S_t$.
\item Calibration condition: unchanged, but also stated under one measure, (\ref{eq: singleMeasure}). The second system of particles under $\mathbb{Q}^{f^T}$, and with it the section on exact integration under $\mathbb{Q}^{f^T}$, are removed.
\item Change of measure $\mathbb{Q}^d\to\mathbb{Q}^f$: the spot driver is shifted by $\sigma(t,S_t)$, not by $\rho_{Sd}\sigma(t,S_t)$.
\item Foreign rate: centred at the market forward rate. The terms in $\mathbb{E}[\sigma(s,S_s)]$ disappear, as does a drift term that had been computed with $\kappa_d$ in place of $\kappa_f$.
\item Discretised forward: sign of the drift, $I^{2,d}_j$ in the domestic variance term, and the coefficient $K^{df}_j$.
\item General Malliavin formula (\ref{eq: threeFactorGeneralMalliavin}): the foreign term has the factor $2/K$ (it was $2S_0/K^2$).
\item The lemmas on $\MD_sr^d_t$, $\MD_sr^f_t$ and $\MD_sS_t$, including the ``second approach'', are replaced by Propositions \ref{prop: integratedSensitivities} and \ref{prop: main}.
\item Sections \ref{subsec: checks} and \ref{subsec: caveats} are new. Bibliography: duplicate entry removed, seven references added.
\end{itemize}

\begin{thebibliography}{99}

\bibitem{guyonLabordere}
  Julien Guyon, Pierre Henry-Labord\`ere,
  \emph{Nonlinear Option Pricing},
  Chapman \& Hall/CRC, 2013.

\bibitem{deelstraRayee}
  Griselda Deelstra, Gr\'egory Ray\'ee,
  \emph{Local Volatility Pricing Models for Long-dated FX Derivatives},
  Applied Mathematical Finance 20(4), 380--402, 2013.

\bibitem{clark}
  Iain J. Clark,
  \emph{Foreign Exchange Option Pricing},
  2011.

\bibitem{nualart}
  David Nualart,
  \emph{The Malliavin Calculus and Related Topics},
  2005.

\bibitem{ewald}
  Christian-Oliver Ewald,
  \emph{A note on the Malliavin derivative operator under change of variable},
  2009.

\bibitem{hybridModel}
  Lech Aleksander Grzelak,
  \emph{Equity and Foreign Exchange Hybrid Models for Pricing Long-Maturity Financial Derivatives},
  2009.

\bibitem{bruti}
  Eckhard Platen, Nicola Bruti-Liberati,
  \emph{Numerical Solution of Stochastic Differential Equations with Jumps in Finance},
  2010.

\bibitem{ogetbilGanesanHientzsch}
  Orcan \"Ogetbil, Narayan Ganesan, Bernhard Hientzsch,
  \emph{Calibrating Local Volatility Models with Stochastic Drift and Diffusion},
  International Journal of Theoretical and Applied Finance 25(2), 2022. arXiv:2009.14764.

\bibitem{ogetbilHientzsch}
  Orcan \"Ogetbil, Bernhard Hientzsch,
  \emph{Extensions of Dupire Formula: Stochastic Interest Rates and Stochastic Local Volatility},
  arXiv:2005.05530, 2020.

\bibitem{cozma}
  Andrei Cozma, Matthieu Mariapragassam, Christoph Reisinger,
  \emph{Calibration of a Hybrid Local-Stochastic Volatility Stochastic Rates Model with a Control Variate Particle Method},
  SIAM Journal on Financial Mathematics 10(1), 181--213, 2019.

\bibitem{josephLoeperObloj}
  Benjamin Joseph, Gr\'egoire Loeper, Jan Ob{\l}\'oj,
  \emph{Calibration of Local Volatility Models with Stochastic Interest Rates using Optimal Transport},
  Finance and Stochastics 30, 397--439, 2026.

\bibitem{abergelTachet}
  Fr\'ed\'eric Abergel, R\'emi Tachet,
  \emph{A nonlinear partial integro-differential equation from mathematical finance},
  Discrete and Continuous Dynamical Systems - Series A 27(3), 907--917, 2010.

\bibitem{jourdainZhou}
  Benjamin Jourdain, Alexandre Zhou,
  \emph{Existence of a calibrated regime switching local volatility model},
  Mathematical Finance 30(2), 501--546, 2020.

\bibitem{lackerShkolnikovZhang}
  Daniel Lacker, Mykhaylo Shkolnikov, Jiacheng Zhang,
  \emph{Inverting the Markovian projection, with an application to local stochastic volatility models},
  Annals of Probability 48(5), 2189--2211, 2020.

\end{thebibliography}

\end{document}
